\section{引言}
现代空域场景日益呈现多飞行器协同交会特征：两枚或更多非合作飞行器（non-cooperative vehicle, NCV）沿刻意错开的交会几何与剩余飞行时间逼近同一指定飞行器\cite{yan2020b,shalumov2019}。对规避者而言，这使规避问题发生了质的变化：增大与其中一枚NCV终端间隔的机动，可能同时缩小与另一枚的间隔；且每一次机动都会消耗随后用于到达指定目标的控制能力。因此，设计一种在同时面对两枚NCV时仍保持有效、并兼顾终端精度的规避策略，已成为制导设计的核心问题。

现有规避制导研究大致沿两条主线展开。第一条主线是在改进的制导结构中嵌入规避需求，如面向采用改进比例导引（PN）NCV的最优规避律\cite{feng2023}、适用于复杂多飞行器场景的有限时间收敛制导律（FTCG）\cite{ding2024}、用于三体交会的协同鲁棒视线（LOS）制导律\cite{luo2023}，以及逃逸距离可控的规避律\cite{yu2023}。第二条主线是将交会过程建模为规避博弈，由博弈最优性条件导出规避者策略\cite{turetsky2003}，并进一步推广到三方高超声速交会\cite{yan2020a}、多方微分对策\cite{liang2022}以及安全关键框架下的协同规避\cite{chen2024}。Shaferman的研究进一步表明，针对采用现代线性制导律的NCV，可以显式构造近最优规避策略\cite{shaferman2021}，因此可将对方的制导结构视为已知先验信息。然而，上述研究大多以终端间隔为主导目标，甚至将其作为唯一目标，所得制导指令往往使可用法向过载趋于饱和，导致后续的目标到达任务难以完成。

能量感知规避因此受到关注。Li等人针对高超声速变体飞行器提出了节能机动策略\cite{li2024cja}；Zhuang等人将深度强化学习（DRL）应用于六自由度模型，以在姿态稳定性与能耗之间权衡\cite{zhuang2024}；Yao等人则在FTCG律中通过精确时序控制优化能量\cite{yao2024}。Wang等人将终端间隔驱动至给定临界安全值而非最大化，从而限制不必要的控制消耗\cite{wang2024}。规避与继续原任务之间的耦合，已由Qi等人针对受保护目标\cite{qi2017}以及Ren等人通过分层规避—到达策略\cite{ren2025}加以研究。这些工作表明最大规避与终端精度是相互冲突的需求，但其要么只处理单枚NCV，要么将两个子任务分层解耦，因而两者的交互从未被纳入同一最优性条件。

明确考虑两枚或更多NCV的研究仍然稀少，并在维数、解的形式或二者上受到限制。Ben-Asher等人在路径角约束下求解了针对两枚NCV的最优规避，但是平面设定且为开环解\cite{benasher1988}。Shen等人优化了高超声速滑翔飞行器遭遇两枚NCV时的规避轨迹，同样为离线求解，因而无法响应飞行中的几何扰动\cite{shen2022}。Zhang等人针对多枚轨道NCV构建了逃逸区，这一几何方法难以适用于剩余飞行时间紧张的大气层末段交会\cite{zhang2023}。Hao等人确实用强化学习研究了双飞行器问题，但其参数化策略建立在二维运动学模型上，无法刻画支配真实三维（3D）交会的面外耦合\cite{hao2025}。目前仍缺乏一种闭环解析制导律，能将规避两枚NCV、到达目标与控制能耗作为三维空间中的同一耦合优化问题来处理。零控偏差量（zero-effort miss, ZEM）可将线性化交会的终端性能压缩为单一标量\cite{zarchan2019}，并已被证明是反馈制导\cite{guo2013}与逆最优制导\cite{ma2022}的有效综合变量，从而为同时缩减多次交会提供了自然途径。

即便已有解析制导律，其权重系数也无法先验固定：两项规避项、到达项与能量项的相对重要性取决于交会几何，因而权值选取等价于必须沿轨迹求解的非线性优化。深度强化学习适合承担这一任务，已被用于由纯角度测量学习末制导\cite{gaudet2020}、实现无需在线数值优化的计算制导\cite{he2021}，以及处理存在干扰目标时的末制导\cite{deng2023}。确定性策略梯度方法适合此类连续参数调度\cite{lillicrap2016}，而其双延迟变体（TD3）可缓解使朴素形式在多约束下不可靠的价值过估计与策略振荡\cite{fujimoto2018}。然而，从零开始学习样本效率低且易陷入较差局部最优，因此训练中常注入示范数据——通过预填充经验回放缓冲\cite{hester2018,vecerik2017}、将策略更新正则化至示范动作\cite{fujimoto2021}，或对交会数据进行对抗模仿\cite{hu2022,hui2025}。但在现有制导应用中，这类示范多来自启发式规则或其他已学策略，而非底层最优控制问题，其所携带的先验本身并非最优。

另一缺口在于测量的真实性。零控偏差量动力学的系数是视线角与角速度的函数，若直接由含噪机载传感器输出计算，将产生有偏且相互相关的误差并直接传入制导指令，因此在线估计必不可少\cite{liu2025}。带遗忘因子的递推最小二乘（RLS）是跟踪此类时变系数的标准工具\cite{paleologu2008,bruce2020}，但在某些几何下投影系数激励较弱，无正则化的递推会发生漂移。附加$\ell_1$惩罚可抑制虚假分量并保持递推形式\cite{angelosante2010}，其思路与控制方程的稀疏辨识相一致\cite{brunton2016}。

为弥补上述不足，本文提出了一种面向双NCV的智能三维规避算法。主要贡献如下。

\begin{enumerate}
\item 建立了三维四体相对运动学模型，将规避与目标到达一体化问题由单NCV推广到双NCV，并从原有研究只考虑二维运动学拓展到三维。
\item 设计了同时考虑规避与目标到达的一体化最优制导律，防止过度规避导致目标到达失败。利用极大值原理推出的三维最优制导律充分考虑了规避两枚NCV与到达目标等多个任务间存在的耦合，并通过求解$3\times 3$耦合零控偏差量系统得到闭环解析形式，生成指令无需在线数值优化。
\item 考虑实际应用中存在的测量误差，设计了一种带稀疏正则化与遗忘因子的最小二乘辨识算法，对状态方程矩阵进行结构化辨识，并与制导律闭环，使辨识量替代其解析计算值。
\item 利用打靶法生成考虑各类初始条件的专家轨迹，设计了BC-TD3框架，使用专家轨迹进行预填充后开启强化学习训练，提高多约束复杂任务下的收敛速度与泛化性。在300余次蒙特卡洛仿真中，所提方法在无测量噪声时的双NCV规避成功率与目标到达成功率分别为94.7\%和94.3\%，结合在线辨识后在有噪声条件下分别达到96.7\%和96.0\%。
\end{enumerate}

本文其余部分安排如下。第2节建立三维四体相对运动学并给出一体化优化问题。第3节推导一体化最优制导律，第4节给出结构化稀疏遗忘因子递推最小二乘（FF-RLS）辨识方案。第5节介绍专家轨迹生成与BC-TD3框架。第6节给出蒙特卡洛仿真结果，第7节总结全文。
